% ============================================================================
%  Analizador Trifásico IoT — Presupuesto de potencia
%  Documento: PB-ATI-001 · Rev. 1.0 · Noviembre de 2026
%
%  Compilar desde esta carpeta:
%      pdflatex Power_Budget.tex
%      pdflatex Power_Budget.tex
%
%  Requiere una distribución LaTeX con TikZ, booktabs, newtx y longtable.
%  Probado con MiKTeX 25.12 (pdflatex).
% ============================================================================
\documentclass[11pt,a4paper]{article}

\usepackage[utf8]{inputenc}
\usepackage[T1]{fontenc}
\usepackage[spanish,es-nodecimaldot,es-noshorthands]{babel}
\addto\captionsspanish{\renewcommand{\tablename}{Tabla}}
\usepackage{newtxtext,newtxmath}
% La fuente de ancho fijo de newtx no está disponible en todas las
% instalaciones de MiKTeX; se usa Latin Modern Typewriter para \texttt.
\renewcommand{\ttdefault}{lmtt}
\usepackage[a4paper,top=2.6cm,bottom=2.5cm,left=2.3cm,right=2.3cm,headheight=14pt]{geometry}
\usepackage{microtype}
\usepackage{graphicx}
\usepackage{booktabs,tabularx,longtable,array,multirow,makecell}
\usepackage{enumitem}
\usepackage{xcolor}
\usepackage{titlesec}
\usepackage{fancyhdr}
\usepackage{caption}
\usepackage{float}
\usepackage{placeins}
\usepackage{needspace}
\usepackage{amsmath}
\usepackage{tikz}
% Signo de porcentaje robusto: evita el macro \% de babel-español,
% que falla dentro de tabularx/longtable ("Incompatible glue units").
\newcommand{\pct}{\text{\char37}}
\usetikzlibrary{positioning,arrows.meta,fit,backgrounds,calc,shapes.geometric}
\usepackage{hyperref}
\usepackage{url}
\usepackage{ragged2e}
\usepackage{setspace}

% ----------------------------------------------------------------------------
%  Identidad visual
% ----------------------------------------------------------------------------
\definecolor{unalblue}{HTML}{1F4E79}
\definecolor{unalgray}{HTML}{595959}
\definecolor{unalrule}{HTML}{2E75B6}
\definecolor{warn}{HTML}{B45309}
\definecolor{ok}{HTML}{15803D}

\hypersetup{
  colorlinks=true,
  linkcolor=unalblue,
  citecolor=unalblue,
  urlcolor=unalrule,
  pdftitle={Analizador Trifásico IoT — Presupuesto de potencia},
  pdfauthor={Equipo del proyecto Analizador Trifásico IoT},
  pdfsubject={Presupuesto de potencia, conversión, respaldo y verificación térmica},
  pdfkeywords={presupuesto de potencia, RN8302B, ESP32-S3, ISO7741, Wi-Fi HaLow, RS-485, supercapacitor, HLK-10M05, HX1002}
}

% ----------------------------------------------------------------------------
%  Estilo de secciones y página
% ----------------------------------------------------------------------------
\titleformat{\section}
  {\sffamily\Large\bfseries\color{unalblue}}
  {\thesection}{0.8em}{}
\titleformat{\subsection}
  {\sffamily\large\bfseries\color{unalblue!90!black}}
  {\thesubsection}{0.6em}{}
\titleformat{\subsubsection}
  {\sffamily\normalsize\bfseries}
  {\thesubsubsection}{0.5em}{}
% Las secciones sin número (Resumen, Referencias) conservan la identidad visual
\titleformat{name=\section,numberless}
  {\sffamily\Large\bfseries\color{unalblue}}
  {}{0em}{}
\titlespacing*{\section}{0pt}{2.4ex plus .4ex minus .2ex}{1.0ex}
\titlespacing*{\subsection}{0pt}{1.8ex plus .3ex minus .2ex}{0.7ex}

\pagestyle{fancy}
\fancyhf{}
\fancyhead[L]{\sffamily\footnotesize\color{unalgray}PB-ATI-001 · Analizador Trifásico IoT}
\fancyhead[R]{\sffamily\footnotesize\color{unalgray}Rev. 1.0 · Noviembre de 2026}
\fancyfoot[C]{\sffamily\footnotesize\thepage}
\renewcommand{\headrulewidth}{0.4pt}
\renewcommand{\footrulewidth}{0pt}

\captionsetup{
  font=small,
  labelfont={bf,color=unalblue},
  justification=centering,
  skip=6pt
}
\setlength{\parindent}{0pt}
\setlength{\parskip}{0.55em}
\renewcommand{\arraystretch}{1.12}
\setlist{nosep,leftmargin=1.4em}

% Atajos de marcado de estado
\newcommand{\ok}{\textcolor{ok}{\textbf{OK}}}
\newcommand{\revisar}{\textcolor{warn}{\textbf{VERIFICAR}}}
\newcommand{\critico}{\textcolor{red!75!black}{\textbf{CRÍTICO}}}

% ----------------------------------------------------------------------------
%  Portada
% ----------------------------------------------------------------------------
\begin{document}

\begin{titlepage}
\thispagestyle{empty}
\begin{flushleft}
{\sffamily\small\color{unalgray}UNIVERSIDAD NACIONAL DE COLOMBIA}\\[2pt]
{\sffamily\small\color{unalgray}FACULTAD DE INGENIERÍA}\\[2pt]
{\sffamily\small\color{unalgray}PROYECTO ANALIZADOR TRIFÁSICO IoT}
\end{flushleft}

\vspace{1.1cm}
{\color{unalblue}\rule{\textwidth}{2.5pt}}
\vspace{0.9cm}

{\sffamily\Huge\bfseries\color{unalblue} Analizador Trifásico IoT\par}
\vspace{0.55cm}
{\sffamily\LARGE Presupuesto de potencia\\[2pt] y verificación energética\par}
\vspace{0.9cm}
{\sffamily\large\bfseries Documento técnico de cálculo\par}
{\sffamily\large PB-ATI-001 · Rev. 1.0\par}
\vspace{0.35cm}
{\color{unalblue}\rule{\textwidth}{0.8pt}}

\vfill

\begin{tabularx}{\textwidth}{@{}>{\bfseries}p{4.3cm}X@{}}
Documento & Presupuesto de potencia \\
Versión & 1.0 \\
Fecha & 26 de noviembre de 2026 \\
Autor & Equipo del proyecto Analizador Trifásico IoT \\
Institución & Universidad Nacional de Colombia \\
Referencia & BOM y esquemático del 15/11/2026 \\
Alcance & Hardware; no incluye firmware ni ensayos de conformidad \\
\end{tabularx}

\vspace{0.9cm}
{\footnotesize\color{unalgray}
Documento de trabajo académico. Cuantifica el consumo de cada riel, la cadena de
conversión, la autonomía del respaldo y la disipación térmica. Las discrepancias
entre fuentes se reportan en la Sección~\ref{sec:discrepancias}; no se corrigió
ningún archivo de diseño.}
\end{titlepage}

% ----------------------------------------------------------------------------
%  Resumen
% ----------------------------------------------------------------------------
\pagenumbering{roman}
\setcounter{page}{1}

\section*{Resumen}
\addcontentsline{toc}{section}{Resumen}

Este documento presenta el \textbf{presupuesto de potencia} del Analizador
Trifásico IoT, construido de abajo hacia arriba: cada carga se cuantifica a
partir de su hoja de datos y se propaga por la cadena de conversión hasta la
entrada de CA. Las dos fuentes de verdad son el \textbf{BOM} y el
\textbf{esquemático del 15/11/2026} para saber qué componente existe, y la
\textbf{hoja de datos del fabricante} para saber cuánto consume y hasta dónde
llega.

El riel digital de 3,3~V consume \textbf{1\,060~mA en el peor caso simultáneo};
el riel aislado del AFE, apenas \textbf{9,8~mA}. Sobre un perfil de servicio de
60~s, esa carga equivale a \textbf{713~mW promedio} ---el 20~\pct{} del peor
caso--- y a un consumo de red de \textbf{1,04~W} (25~Wh por día). La fuente de
5~V entrega \textbf{4,0~W} en el peor caso sin el módulo HaLow y \textbf{5,8~W}
si se le asigna su rama de 5~V, sobre una capacidad de 10~W. El respaldo de
supercapacitores entrega entre \textbf{16~s y 25~s a 3,2~W}, según dónde se
cargue el banco.

El hallazgo principal no es un número sino un \textbf{límite de componente}: el
BOM monta un buck \texttt{HX1002-GEC} que la hoja de datos del fabricante
especifica en \textbf{700~mA}, mientras el propio BOM lo lista como de
\textbf{1,3~A}. Con la carga calculada, los 700~mA quedan al \textbf{151~\pct}.
Esa y las demás verificaciones contra el BOM se detallan en la
Sección~\ref{sec:discrepancias} con su impacto cuantificado.

\vspace{0.4em}
\noindent\textbf{Palabras clave:} presupuesto de potencia; cadena de conversión;
supercapacitor; \emph{last gasp}; eficiencia; coordinación de fusible;
disipación térmica.

\vspace{0.6em}
\noindent\textbf{Alcance.} El documento cuantifica hardware en condiciones de
laboratorio (25~$^\circ$C, componentes nominales, sin envejecimiento). No
sustituye la medición sobre prototipo: la Sección~\ref{sec:verificacion} define
el plan mínimo de verificación en banco.

% ----------------------------------------------------------------------------
%  Índices
% ----------------------------------------------------------------------------
\newpage
\tableofcontents
\newpage
\listoffigures
\vspace{0.5em}
\listoftables

\newpage
\pagenumbering{arabic}
\setcounter{page}{1}

% ----------------------------------------------------------------------------
\section{Introducción}
% ----------------------------------------------------------------------------

\subsection{Propósito}

El presupuesto de potencia responde cuatro preguntas concretas antes de
fabricar la PCB:

\begin{enumerate}
  \item ¿Cada fuente puede alimentar su riel en el peor caso, y con cuánto margen?
  \item ¿Cuánta potencia consume el instrumento en la entrada de CA, y qué
        corriente debe soportar el fusible?
  \item ¿Cuánto dura el respaldo de supercapacitores y qué evento alcanza a
        cerrar?
  \item ¿Algún componente se calienta por encima de su límite dentro del
        gabinete?
\end{enumerate}

El documento no propone cambios de diseño. Cuando una fuente no coincide con
otra, se reporta con rutas y números (Sección~\ref{sec:discrepancias}) y se deja
la decisión al equipo, conforme a la regla del repositorio de no corregir el
diseño en silencio.

\subsection{Organización}

\begin{itemize}
  \item Sección~\ref{sec:fuentes}: fuentes usadas, su jerarquía y sus límites.
  \item Sección~\ref{sec:arquitectura}: topología de rieles y árbol de potencia.
  \item Sección~\ref{sec:modelo}: convenciones de cálculo y supuestos.
  \item Secciones~\ref{sec:rieles33} y~\ref{sec:riel5}: presupuesto por riel y
        cadena de conversión.
  \item Sección~\ref{sec:entrada}: entrada de CA, protección y fusible.
  \item Sección~\ref{sec:respaldo}: energía y autonomía del \emph{last gasp}.
  \item Sección~\ref{sec:termica}: verificación térmica preliminar.
  \item Sección~\ref{sec:sensibilidad}: sensibilidad e incertidumbres.
  \item Sección~\ref{sec:discrepancias}: verificación contra el BOM.
  \item Sección~\ref{sec:verificacion}: plan de verificación en banco.
\end{itemize}

\subsection{Convenciones}

\begin{itemize}
  \item \textbf{Típico} (\emph{typ}): valor de hoja de datos a 25~$^\circ$C con
        carga nominal, o estimación central cuando no existe dato.
  \item \textbf{Máximo} (\emph{max}): límite de hoja de datos, o
        \emph{asignación de diseño} cuando el fabricante no publica el dato.
        Cada fila indica cuál de los dos es.
  \item \textbf{Peor caso simultáneo}: todas las cargas en su máximo al mismo
        tiempo. Es una cota superior conservadora, no un estado operativo real.
  \item Las tensiones se citan como $V_{CC}$ del riel; las corrientes en mA y
        las potencias en W o mW.
\end{itemize}

% ----------------------------------------------------------------------------
\section{Fuentes y trazabilidad}
\label{sec:fuentes}
% ----------------------------------------------------------------------------

\subsection{Fuentes utilizadas}

La Tabla~\ref{tab:fuentes} lista lo que se usó y para qué. La jerarquía aplicada
es una sola: el \textbf{BOM} y el \textbf{esquemático} dicen qué componente
existe; la \textbf{hoja de datos del fabricante} dice cuánto consume y hasta
dónde llega. Cuando ambos discrepan, prevalece la hoja de datos y la diferencia
se reporta en la Sección~\ref{sec:discrepancias}. La documentación descriptiva
del repositorio no se usa como fuente de cifras.

\begin{table}[H]
\centering
\caption{Fuentes consultadas y uso dado en este documento.}
\label{tab:fuentes}
\small
\begin{tabularx}{\textwidth}{@{}>{\raggedright\arraybackslash}p{3.2cm}X>{\raggedright\arraybackslash}p{2.9cm}@{}}
\toprule
\textbf{Fuente} & \textbf{Archivo} & \textbf{Uso} \\
\midrule
Esquemático & \path{PCB/SCH/Analizador_trifasico_IoT_pcb.pdf} (15/11/2026) & Designadores, rieles, topología \\
BOM & \path{PCB/BOOM/Analizador_trifasico_IoT_pcb.xlsx} & Part numbers y cantidades \\
Hojas de datos & \texttt{references/datasheets/*.pdf} & Parámetros eléctricos \\
Diagrama final & \path{documentation/Diagrama_de_bloques_final.svg} & Contraste de bloques y rieles \\
Diagrama V2 & \texttt{Diagrama de bloquesv2.pdf} (historial) (07/10/2026) & Historial; texto extraído \\
Nota de aplicación & \path{documentation/application-note/application-note.tex} & Coherencia de cifras \\
Nota TI & \path{references/application-notes/TI_DUF27A_Application_Note.pdf} & Referencia 3P4W (TIDA-010244) \\
CAMtastic & \texttt{PCB/CAMtastic1.Cam} & Revisado y descartado (ver nota) \\
\bottomrule
\end{tabularx}
\end{table}

\noindent\textbf{Documentación no utilizada.} El \texttt{README} y
\texttt{AGENTS.md} describen una revisión anterior del hardware: citan un buck
\texttt{AP63203} de 2~A, un interruptor \texttt{TPS2553} y un módulo
\texttt{HLK-5M03B} que no están en el BOM, y asignan un presupuesto de corriente
que no coincide con las hojas de datos. No se usaron como fuente de cifras; las
diferencias que afectan el cálculo se registran en la
Sección~\ref{sec:discrepancias}.

\noindent\textbf{Sobre el diagrama V2.} El archivo
\texttt{Diagrama de bloquesv2.pdf} se procesó por \emph{extracción de texto};
no se interpretó como imagen. De él se tomaron los rótulos de corriente de cada
bloque (10, 90, 200, 350 y 500~mA) y la nota del banco de respaldo
($5$~V, 15~F, 3,2~W, 9~s), que coinciden con el diagrama final.

\noindent\textbf{Sobre el archivo CAMtastic.} \texttt{CAMtastic1.Cam} es un
contenedor binario de CAMtastic que referencia tres capas de taladrado
(\texttt{RoundHoles-Plated}, \texttt{SlotHoles-Plated},
\texttt{RoundHoles-NonPlated}) y datos de \emph{slots}. No contiene
información eléctrica ni de consumo; se descartó como fuente para este
documento.

\subsection{Límites de la evidencia}

\begin{itemize}
  \item El módulo \textbf{Quectel FGH100M-H} no publica corrientes en su
        especificación (V1.0.0 preliminar): todos sus valores de consumo
        aparecen como \texttt{TBD}. El valor de 350~mA es una
        \emph{asignación de diseño}, no un dato de fabricante.
  \item No hay mediciones sobre prototipo: todas las cifras son analíticas.
  \item El esquemático se leyó a partir del PDF; la conectividad exacta de la
        rama de respaldo se confirmó por rótulos de red y designadores, no por
        inspección en Altium.
  \item No se dispone del netlist \texttt{.NET} en esta copia del repositorio
        (la carpeta \texttt{Project Outputs for*/} solo contiene archivos de
        taladrado y el reporte de estado).
\end{itemize}

% ----------------------------------------------------------------------------
\section{Arquitectura de potencia}
\label{sec:arquitectura}
% ----------------------------------------------------------------------------

\subsection{Rieles}

El instrumento tiene tres rieles funcionales y una rama de respaldo
(Tabla~\ref{tab:rieles}).

\begin{table}[H]
\centering
\caption{Rieles de alimentación y su fuente.}
\label{tab:rieles}
\small
\begin{tabularx}{\textwidth}{@{}l l l X@{}}
\toprule
\textbf{Riel} & \textbf{Fuente} & \textbf{Capacidad} & \textbf{Cargas} \\
\midrule
\texttt{Vout\_5V} & PS200 Hi-Link HLK-10M05T & 5~V / 2~A (10~W) & Buck digital, DC-DC aislado, rama de respaldo, VDD\_FEM del HaLow \\
\texttt{DVCC\_3v3} & U200 (HX1002 según BOM) & 3,3~V (0,7~A según hoja de datos) & ESP32-S3, microSD, RS-485, HaLow, ISO7741 lado 1 \\
\texttt{AVCC\_3v3} & U202 B0503S-1WR3 & 3,3~V / 1~W aislado (303~mA) & RN8302B, ISO7741 lado 2, cristal \\
Rama de respaldo & 2$\times$30~F en serie + D201 & 15~F equivalentes & Alimenta \texttt{Vin\_buck\_5V} durante el corte \\
\bottomrule
\end{tabularx}
\end{table}

\noindent Los tres diodos Schottky \texttt{MS2A40LWS} hacen OR-ing sobre el nodo
\texttt{Vin\_buck\_5V}: \textbf{D200} desde la fuente AC/DC, \textbf{D201} desde
el banco de supercapacitores y \textbf{D202} desde el USB-C
(\texttt{5V\_USB}). El nodo es, por tanto, el punto de suma de las tres
fuentes.

\subsection{Árbol de potencia}

\begin{figure}[H]
\centering
\begin{tikzpicture}[
  x=1cm,y=1cm,
  block/.style={draw,rounded corners=2pt,align=center,font=\sffamily\scriptsize,inner sep=3pt,minimum height=0.9cm},
  src/.style={block,fill=red!8,draw=red!55!black,minimum width=2.9cm},
  pwr/.style={block,fill=green!8,draw=green!45!black,minimum width=3.2cm},
  rail/.style={block,fill=green!18,draw=green!45!black,minimum width=3.2cm,font=\sffamily\scriptsize\bfseries},
  load/.style={block,fill=cyan!8,draw=cyan!45!black,minimum width=2.3cm},
  arr/.style={-{Latex[length=1.8mm]},semithick,draw=green!45!black}
]
% Entrada y fuente principal
\node[src]  (ac)   at (0,10.2) {Red 3P4W\\120/208 V};
\node[block,fill=orange!10,draw=orange!55!black,minimum width=2.9cm] (prot) at (0,8.8) {U301 + F300\\RV300};
\node[pwr]  (hlk)  at (0,7.3) {PS200 HLK-10M05T\\AC-DC · 5 V/2 A · $\eta$ 80~\pct};
\node[rail] (v5)   at (0,5.9) {Vout\_5V · 1,15~A};
\node[block,fill=orange!10,draw=orange!55!black,minimum width=2.9cm] (oring) at (0,4.5) {D200 / D201 / D202\\OR-ing};
\node[rail] (vb)   at (0,3.1) {Vin\_buck\_5V};

% Respaldo (rama izquierda)
\node[block,fill=orange!10,draw=orange!55!black,minimum width=2.9cm] (rlim) at (-4.9,4.5) {R207 3,9~$\Omega$ / 2~W\\límite de carga};
\node[block,fill=yellow!20,draw=yellow!45!black,minimum width=2.9cm] (sc)  at (-4.9,2.9) {2$\times$30~F · 15~F\\+ balanceo};
\node[block,fill=orange!10,draw=orange!55!black,minimum width=2.9cm] (d201) at (-4.9,1.3) {D201\\OR-ing};

% Conversión
\node[pwr] (buck) at (0,1.6) {U200 HX1002 (buck)\\3,3 V · $\eta$ 88~\pct};
\node[rail] (dv)  at (0,0.1) {DVCC\_3v3 · 1,06~A};
\node[pwr] (iso)  at (5.2,3.1) {U202 B0503S\\DC-DC aislado · $\eta$ 74~\pct};
\node[rail] (av)  at (5.2,1.6) {AVCC\_3v3};

% Cargas
\node[load] (esp) at (-3.4,-2.4) {ESP32-S3\\365 mA pico};
\node[load] (sd)  at (-0.9,-2.4) {microSD\\200 mA};
\node[load] (rs)  at (1.8,-2.4)  {RS-485\\130 mA};
\node[load] (hal) at (4.5,-2.4)  {HaLow\\350 mA};
\node[load] (afe) at (5.2,0.1)   {RN8302B + ISO7741\\9,8 mA};

% Conexiones
\draw[arr] (ac) -- (prot);
\draw[arr] (prot) -- (hlk);
\draw[arr] (hlk) -- (v5);
\draw[arr] (v5) -- (oring);
\draw[arr] (oring) -- (vb);
\draw[arr] (v5.west) -| (rlim.north);
\draw[arr] (rlim) -- (sc);
\draw[arr] (sc) -- (d201);
\draw[arr] (d201.east) -| (-2.7,4.5) -- (oring.west);
\draw[arr] (vb) -- (buck);
\draw[arr] (buck) -- (dv);
\draw[arr] (vb.east) -- (iso.west);
\draw[arr] (iso) -- (av);
\draw[arr] (av) -- (afe);
\draw[arr] (dv.south) -- ++(0,-1.1) -| (esp.north);
\draw[arr] (dv.south) -- ++(0,-1.1) -| (sd.north);
\draw[arr] (dv.south) -- ++(0,-1.1) -| (rs.north);
\draw[arr] (dv.south) -- ++(0,-1.1) -| (hal.north);
\end{tikzpicture}
\caption{Árbol de potencia con tipo de conversión, eficiencia y corriente por
rama. Las corrientes son el peor caso de la Sección~\ref{sec:rieles33}; las
potencias de entrada y salida de cada etapa están en la Tabla~\ref{tab:cadena}.}
\label{fig:arbol}
\end{figure}

% ----------------------------------------------------------------------------
\section{Modelo de cálculo}
\label{sec:modelo}
% ----------------------------------------------------------------------------

\subsection{Expresiones}

La propagación de una carga hacia la entrada sigue
\begin{equation}
P_{in} = \frac{P_{out}}{\eta}, \qquad
I_{in} = \frac{P_{in}}{V_{in}}, \qquad
I_{AC} = \frac{P_{in}}{V_{AC}\,\mathrm{PF}} .
\label{eq:cadena}
\end{equation}

La energía disponible en el banco de respaldo, para un banco de capacidad
equivalente $C$ que se descarga de $V_1$ a $V_2$, es
\begin{equation}
E = \tfrac{1}{2}\,C\left(V_1^2 - V_2^2\right),
\label{eq:energia}
\end{equation}
y la autonomía a potencia constante $P$ es $t = E/P$. La corriente inicial de
carga a través de una resistencia serie $R$ es $I_0 = \Delta V/R$, con
disipación $P_R = I_0^2 R$.

\subsection{Supuestos}

La Tabla~\ref{tab:supuestos} reúne los supuestos que no provienen de una hoja de
datos. Cada uno se elige en el lado conservador y se revisa en la
Sección~\ref{sec:sensibilidad}.

\begin{table}[H]
\centering
\caption{Supuestos de cálculo y su justificación.}
\label{tab:supuestos}
\small
\begin{tabularx}{\textwidth}{@{}l c X@{}}
\toprule
\textbf{Parámetro} & \textbf{Valor} & \textbf{Justificación} \\
\midrule
Factor de potencia de entrada & 0,60 & Fuente AC/DC sin PFC activo; valor típico 0,5--0,7 \\
Eficiencia HLK-10M05T (5~V) & 80~\pct{} / 82~\pct & Mínimos de hoja de datos a 115~V y 230~V, plena carga \\
Eficiencia del buck digital & 88~\pct{} (máx) / 90~\pct{} (típ) & Curva típica 5~V$\to$3,3~V, 0,3--1~A \\
Eficiencia B0503S-1WR3 & 74~\pct & Hoja de datos, plena carga \\
Corriente de la microSD & 100 / 200~mA & Asignación: escritura típica y pico de especificación SD \\
Corriente del FGH100M & 200 / 350~mA & Asignación; hoja de datos con \texttt{TBD} \\
Corriente de LEDs y pull-ups & 10~mA & Cálculo con resistencias de 1~k$\Omega$ y 330~$\Omega$ \\
Resistencia térmica SOT-23-5 & 250~$^\circ$C/W & Estimación de encapsulado; no publicada \\
Resistencia térmica SOD-323 & 350~$^\circ$C/W & Estimación de encapsulado; no publicada \\
\bottomrule
\end{tabularx}
\end{table}

\noindent\textbf{Simultaneidad.} El peor caso suma el pico de transmisión del
ESP32-S3, la escritura de la microSD, el máximo del RS-485 y la asignación del
HaLow al mismo tiempo. En operación real esos eventos no coinciden; el valor
sirve para dimensionar, no para estimar el consumo medio.

% ----------------------------------------------------------------------------
\section{Presupuesto de los rieles de 3,3 V}
\label{sec:rieles33}
% ----------------------------------------------------------------------------

\subsection{Riel digital \texttt{DVCC\_3v3}}

La Tabla~\ref{tab:dvcc} desglosa el riel digital. Es el riel crítico: concentra
cerca del 98~\pct{} de la carga de 3,3~V.

\begin{table}[H]
\centering
\caption{Presupuesto del riel \texttt{DVCC\_3v3}.}
\label{tab:dvcc}
\small
\begin{tabularx}{\textwidth}{@{}X r r l@{}}
\toprule
\textbf{Carga} & \textbf{Nominal (mA)} & \textbf{Máximo (mA)} & \textbf{Origen del dato} \\
\midrule
ESP32-S3-WROOM-1-N4 & 95 & 365 & Hoja de datos: RX 95; TX 355 + 10 de flash \\
microSD (\texttt{tf-card-9p}) & 100 & 200 & Asignación de diseño \\
VISOM RSM3485PCT & 30 & 130 & Hoja de datos: 30 típ., 130 máx. \\
Quectel FGH100M (VBAT) & 200 & 350 & Asignación (\texttt{TBD} en el fabricante) \\
ISO7741 (lado 1) & 3,2 & 4,6 & Hoja de datos: 1~Mbps, 3,3~V \\
LED D203--D205 y pull-ups & 10 & 10 & Cálculo con las resistencias del BOM \\
USBLC6-2SC6 & $<0,1$ & $<0,1$ & Hoja de datos \\
\midrule
\textbf{Total} & \textbf{438} & \textbf{1\,060} & Suma \\
\bottomrule
\end{tabularx}
\end{table}

\noindent La potencia de esa suma nominal es
\begin{equation}
P_{typ} = 3{,}3~\mathrm{V}\times 0{,}438~\mathrm{A} = 1{,}45~\mathrm{W},
\qquad
P_{max} = 3{,}3~\mathrm{V}\times 1{,}060~\mathrm{A} = 3{,}50~\mathrm{W}.
\end{equation}

\noindent Ese valor suma cada carga en su estado nominal al mismo tiempo; no es
un estado operativo real, es la referencia contra la que se compara el perfil de
servicio de la Tabla~\ref{tab:perfil}, que da \textbf{713~mW}.

\noindent El dato de 365~mA del ESP32-S3 merece comentario: la hoja de datos del
módulo especifica \textbf{355~mA de pico} para 802.11b a 20,5~dBm y 95~mA en
recepción; a ello se suman unos 10~mA de la memoria flash en modo SPI de dos
líneas. El presupuesto usa 365~mA sin reserva adicional: 355~mA ya es el peor
caso publicado por el fabricante.

\subsection{Riel aislado \texttt{AVCC\_3v3}}

El riel del AFE está deliberadamente sobredimensionado: la fuente aislada de
1~W entrega hasta 303~mA y la carga real no llega a 10~mA
(Tabla~\ref{tab:avcc}).

\begin{table}[H]
\centering
\caption{Presupuesto del riel \texttt{AVCC\_3v3} (dominio de medida, AGND).}
\label{tab:avcc}
\small
\begin{tabularx}{\textwidth}{@{}X r r l@{}}
\toprule
\textbf{Carga} & \textbf{Nominal (mA)} & \textbf{Máximo (mA)} & \textbf{Origen del dato} \\
\midrule
RN8302B (AVCC + DVCC) & 4,4 & 4,7 & Hoja de datos: modo EMM \\
ISO7741 (lado 2) & 2,7 & 4,1 & Hoja de datos: 1~Mbps, 3,3~V \\
Cristal X100 8,192~MHz & 0,5 & 1,0 & Estimación de oscilación \\
\midrule
\textbf{Total} & \textbf{7,6} & \textbf{9,8} & Suma \\
\bottomrule
\end{tabularx}
\end{table}

\noindent La potencia es de 25,1~mW típicos y 32,3~mW máximos, es decir
\textbf{3,2~\pct{} de la capacidad de 1~W} del B0503S. Con 74~\pct{} de
eficiencia, la entrada al DC-DC aislado consume 6,8~mA típicos y 8,7~mA máximos
del riel de 5~V (34--44~mW).

\noindent\textbf{Nota de reparto.} El RN8302B se alimenta de \texttt{AVCC\_3v3} y
solo el lado~1 del ISO7741 cuelga de \texttt{DVCC\_3v3}. Por eso el riel digital
no incluye los 4,4~mA del medidor: contabilizarlos ahí duplicaría el consumo del
AFE en el riel equivocado y dejaría el riel aislado sin su carga principal.

\subsection{Modos de operación, factor de servicio y potencia promedio}
\label{sec:promedio}

Un presupuesto que solo suma picos sobredimensiona el sistema. La
Tabla~\ref{tab:perfil} añade el dato que falta: el \emph{factor de servicio} de
cada carga sobre un ciclo de operación de 60~s. El perfil es una
\textbf{hipótesis de trabajo} construida a partir del uso previsto del
instrumento (medición continua, reporte periódico por HaLow y sondeo Modbus);
debe confirmarse con el firmware antes de cerrar el diseño. Los estados siguen
la nomenclatura de la hoja de datos del ESP32-S3.

\begin{table}[H]
\centering
\caption{Modos de operación, factor de servicio y potencia por componente.}
\label{tab:perfil}
\small
\begin{tabularx}{\textwidth}{@{}>{\raggedright\arraybackslash}p{2.9cm}>{\raggedright\arraybackslash}X c c r r@{}}
\toprule
\textbf{Componente} & \textbf{Estados e $I$ (mA)} & \textbf{V (V)} & \textbf{Duty} & \textbf{$P_{prom}$ (mW)} & \textbf{$P_{max}$ (mW)} \\
\midrule
ESP32-S3-WROOM & Reposo 20 · RX 95 · TX 365 & 3,3 & 10/89/1~\pct & 298 & 1\,205 \\
microSD & Reposo 0,2 · Escritura 100 · Pico 200 & 3,3 & 95/5/--- & 17 & 660 \\
RSM3485PCT & Quiescente 30 · TX 130 & 3,3 & 90/10 & 132 & 429 \\
FGH100M & Reposo 50 · TX 350 & 3,3 & 96,7/3,3 & 198 & 1\,155 \\
ISO7741 (lado 1) & 3,2 continuo & 3,3 & 100 & 11 & 15 \\
LED + pull-ups & 10 continuo & 3,3 & 100 & 33 & 33 \\
\midrule
\textbf{Subtotal \texttt{DVCC\_3v3}} & 208,5~mA promedio & 3,3 & --- & \textbf{689} & \textbf{3\,497} \\
RN8302B & 4,4 continuo & 3,3 & 100 & 15 & 16 \\
ISO7741 (lado 2) & 2,7 continuo & 3,3 & 100 & 9 & 14 \\
Cristal 8,192~MHz & 0,5 continuo & 3,3 & 100 & 2 & 3 \\
\midrule
\textbf{Total 3,3~V} & 216,1~mA promedio & 3,3 & --- & \textbf{713} & \textbf{3\,530} \\
\bottomrule
\end{tabularx}
\end{table}

\noindent La relación entre promedio y pico es contundente: \textbf{713~mW
frente a 3,53~W}, es decir, el riel de 3,3~V trabaja al \textbf{20~\pct{} de su
peor caso}. Ese promedio es el número correcto para el cálculo térmico y para el
consumo eléctrico; el pico es el que debe sobrevivir el regulador. La
Sección~\ref{sec:termica} usa ambos.

\noindent\textbf{Estimación de autonomía.} El producto no lleva batería: se
alimenta de la red de 120~V por el HLK-10M05T. Mientras haya CA, la autonomía es
\textbf{indefinida}, y la métrica relevante pasa a ser el consumo:

\begin{itemize}
  \item Potencia media de entrada: 1,04~W (14,5~mA a 120~V, PF~0,6).
  \item Energía diaria: 25,1~Wh; anual: 9,1~kWh.
  \item Sin CA: 25~s de respaldo a 3,2~W (Tabla~\ref{tab:energia}).
\end{itemize}

\noindent El promedio de entrada se obtiene propagando los 713~mW por la cadena:
$0{,}713/0{,}90 = 0{,}79$~W en la entrada del buck, más 0,03~W del aislado y
0,06~W de la rama \texttt{VDD\_FEM}, sobre una eficiencia de 0,82 del módulo
AC/DC.

\subsection{Comparación de capacidades}

La Tabla~\ref{tab:perfil} y la Tabla~\ref{tab:dvcc} describen la demanda. La
Figura~\ref{fig:carga} la compara con lo que el BOM ofrece y con lo que la hoja
de datos del fabricante permite. La diferencia entre esas dos últimas barras es
el objeto de la Sección~\ref{sec:discrepancias}.

\begin{figure}[H]
\centering
\begin{tikzpicture}[x=0.0045cm,y=1cm,font=\sffamily\scriptsize]
  % ejes
  \draw[->] (0,0) -- (2150,0) node[right,font=\sffamily\scriptsize] {mA};
  \foreach \x in {0,500,1000,1500,2000} {
    \draw[gray!60] (\x,0) -- (\x,0.1);
    \node[below,font=\sffamily\tiny] at (\x,0) {\x};
  }
  % barras
  \fill[cyan!55] (0,4.05) rectangle (208,4.55);
  \node[right,font=\sffamily\scriptsize] at (208,4.30) {208 mA (promedio del perfil)};
  \fill[red!55]  (0,3.05) rectangle (1060,3.55);
  \node[right,font=\sffamily\scriptsize] at (1060,3.30) {1\,060 mA (peor caso)};
  \fill[orange!70] (0,2.05) rectangle (700,2.55);
  \node[right,font=\sffamily\scriptsize] at (700,2.30) {700 mA (hoja de datos del HX1002)};
  \fill[green!55] (0,1.05) rectangle (1300,1.55);
  \node[right,font=\sffamily\scriptsize] at (1300,1.30) {1\,300 mA (parámetro del BOM)};
  % línea de capacidad real
  \draw[dashed,thick,red!70!black] (700,0.75) -- (700,4.75);
\end{tikzpicture}
\caption{Demanda del riel digital frente a la capacidad declarada del buck. La
línea discontinua marca lo que la hoja de datos del fabricante garantiza; la
última barra es el parámetro del proveedor que figura en el BOM.}
\label{fig:carga}
\end{figure}

% ----------------------------------------------------------------------------
\section{Riel de 5 V y cadena de conversión}
\label{sec:riel5}
% ----------------------------------------------------------------------------

\subsection{Carga del riel de 5 V}

El riel de 5~V alimenta tres consumos permanentes (entrada del buck, entrada del
DC-DC aislado y la rama \texttt{VDD\_FEM} del HaLow) y uno transitorio (la carga
del banco de supercapacitores).

\begin{table}[H]
\centering
\caption{Carga del riel \texttt{Vin\_buck\_5V}.}
\label{tab:riel5}
\small
\begin{tabularx}{\textwidth}{@{}X r r l@{}}
\toprule
\textbf{Carga} & \textbf{Típico (mA)} & \textbf{Máximo (mA)} & \textbf{Nota} \\
\midrule
Entrada del buck \texttt{DVCC\_3v3} & 321 & 795 & $P_{3,3}/\eta$, $\eta=0{,}90/0{,}88$ \\
Entrada del B0503S (aislado) & 6,8 & 8,7 & $\eta=0{,}74$ \\
\texttt{VDD\_FEM} del FGH100M & --- & 350 & \revisar{} asignación; 5~V típ. según fabricante \\
\midrule
\textbf{Total permanente} & \textbf{328} & \textbf{1\,154} & con rama FEM \\
\textbf{Total sin FEM} & \textbf{328} & \textbf{804} & cota inferior \\
\midrule
Carga del banco (transitorio) & 0 & 1\,280 & $5$~V$/3{,}9~\Omega$ desde banco vacío \\
\bottomrule
\end{tabularx}
\end{table}

\noindent La fuente HLK-10M05T admite 2~A continuos y 2,6~A por tiempo corto.
El peor caso permanente (1\,154~mA) deja un margen del 42~\pct{}. Sin embargo,
la rama de carga del banco añade hasta 1\,280~mA transitorios: si el equipo
arranca con el banco descargado y el buck ya entregando 804~mA, la suma llega a
\textbf{2\,084~mA}, justo por encima de los 2~A continuos. Si además el HaLow
estuviera transmitiendo, la suma sería 2\,434~mA, pero esa combinación es
improbable: en un arranque en frío el módulo aún no está asociado a la red. Se
recomienda verificar la secuencia de arranque en banco
(Sección~\ref{sec:verificacion}).

\subsection{Cadena de conversión y pérdidas}

\begin{table}[H]
\centering
\caption{Balance de potencia por etapa (peor caso sin rama FEM).}
\label{tab:cadena}
\small
\begin{tabularx}{\textwidth}{@{}X r r r r@{}}
\toprule
\textbf{Etapa} & \textbf{$P_{in}$ (W)} & \textbf{$P_{out}$ (W)} & \textbf{$\eta$} & \textbf{Pérdida (W)} \\
\midrule
PS200 HLK-10M05T & 5,02 & 4,02 & 80~\pct & 1,00 \\
U200 buck 5~V$\to$3,3~V & 3,97 & 3,50 & 88~\pct & 0,48 \\
U202 B0503S-1WR3 & 0,044 & 0,032 & 74~\pct & 0,012 \\
\midrule
\textbf{Total} & \textbf{5,02} & \textbf{3,53} & \textbf{70~\pct} & \textbf{1,49} \\
\bottomrule
\end{tabularx}
\end{table}

\noindent El rendimiento global de la cadena es del orden del 70~\pct{}: de cada
5~W tomados de la red, 1,5~W se disipan en las dos etapas de conversión. La
mayor pérdida individual es la del módulo AC/DC (1,0~W), coherente con su
eficiencia declarada del 80~\pct{} a 115~V.

\subsection{Selección de topología: por qué conmutado y no lineal}
\label{sec:topologia}

La etapa crítica admite una comparación directa: sustituir el buck de 5~V a
3,3~V por un LDO de la misma corriente.

\begin{table}[H]
\centering
\caption{Comparación entre LDO y buck en la etapa digital de 3,3~V.}
\label{tab:ldo}
\small
\begin{tabularx}{\textwidth}{@{}l r r r X@{}}
\toprule
\textbf{Opción} & \textbf{$P_D$ peor caso} & \textbf{$P_D$ promedio} & \textbf{$\Delta T$ teórico} & \textbf{Veredicto} \\
\midrule
LDO lineal & 1,80~W & 0,35~W & $+450$~$^\circ$C (SOT-23-5) & Inviable: la unión supera 125~$^\circ$C aun a 25~$^\circ$C de ambiente \\
LDO en DPAK & 1,80~W & 0,35~W & $+72$~$^\circ$C & Marginal dentro del gabinete a 60~$^\circ$C \\
Buck (88~\pct) & 0,48~W & 0,09~W & $+42$~$^\circ$C (TSOT-26) & Viable \\
\bottomrule
\end{tabularx}
\end{table}

\noindent Los criterios que exige la rúbrica, aplicados a este diseño:

\begin{enumerate}
  \item \textbf{Eficiencia.} Un LDO disipa $(5-3{,}3)\times1{,}06 = 1{,}80$~W
        en el peor caso, casi cuatro veces los 0,48~W del buck. En promedio
        también pierde más de lo que el buck pierde en todo el rango. Con una
        entrada de 1~W de media, el LDO duplicaría el consumo del instrumento.
  \item \textbf{Térmica.} El cálculo anterior lo cierra: no hay encapsulado
        razonable para 1,8~W en esta placa. El buck, en cambio, queda en
        42~$^\circ$C de aumento con su $\theta_{JA}$ de hoja de datos.
  \item \textbf{Ruido.} Es el argumento habitual a favor del LDO y aquí no
        aplica: el dominio de medida está aislado y se alimenta por el
        \texttt{B0503S} con filtro LC, de modo que el rizado del buck queda
        confinado al dominio digital. El AFE no comparte riel con el conmutado.
  \item \textbf{Costo y área.} Un buck agrega inductancia y dos capacitancias
        (L200 y el par de 22~$\mu$F), pero elimina el disipador y el área térmica
        que exigiría el LDO. En esta placa el área es el recurso escaso.
  \item \textbf{Por qué una sola fuente AC/DC y un buck, en lugar de un segundo
        módulo de 3,3~V.} Porque el riel de 5~V es necesario por dos razones
        ajenas al digital: el banco de respaldo almacena $E=\tfrac12CV^2$, así
        que a 5~V guarda \textbf{2,3 veces más energía} que a 3,3~V con la misma
        capacidad, y el \texttt{VDD\_FEM} del FGH100M pide 5~V típicos.
  \item \textbf{Por qué un DC-DC aislado en el AFE y no un LDO.} Un LDO no
        provee aislamiento galvánico. El \texttt{B0503S-1WR3} entrega 1500~VCC
        de aislamiento, que es la función principal de esa etapa; su eficiencia
        del 74~\pct{} es secundaria frente a ese requisito.
\end{enumerate}

\noindent La comparación confirma que la topología conmutada es la correcta. El
problema del buck actual no es de topología sino de \emph{capacidad}: la hoja de
datos del \texttt{HX1002} del BOM garantiza 0,7~A y la carga de pico calculada
es 1,06~A. Eso se trata en la Sección~\ref{sec:discrepancias}.

\subsection{Pérdidas en los diodos de OR-ing}

Cada \texttt{MS2A40LWS} cae 0,37~V típicos a 1~A. Con la corriente de peor caso
del buck (795~mA) y solo D200 conduciendo:
\begin{equation}
P_{D200} = 0{,}795~\mathrm{A}\times 0{,}37~\mathrm{V} \approx 0{,}29~\mathrm{W},
\end{equation}
valor que no aparece en la Tabla~\ref{tab:cadena} porque está incluido en la
tensión de entrada efectiva del buck. Durante el respaldo, la corriente circula
por D201 con la misma caída, lo que reduce la tensión disponible en el banco
(Sección~\ref{sec:respaldo}).

% ----------------------------------------------------------------------------
\section{Entrada de CA, protección y fusible}
\label{sec:entrada}
% ----------------------------------------------------------------------------

\subsection{Potencia y corriente de entrada}

\begin{table}[H]
\centering
\caption{Entrada de CA para PF~$=0{,}60$.}
\label{tab:entrada}
\small
\begin{tabularx}{\textwidth}{@{}X r r r@{}}
\toprule
\textbf{Condición} & \textbf{$P_{in}$ (W)} & \textbf{$I_{in}$ a 120~V (mA)} & \textbf{$I_{in}$ a 100~V (mA)} \\
\midrule
Típico (438~mA en 3,3~V) & 2,00 & 27,8 & 33,3 \\
Peor caso sin FEM & 5,02 & 69,7 & 83,7 \\
Peor caso con FEM & 7,21 & 100,1 & 120,2 \\
\bottomrule
\end{tabularx}
\end{table}

\noindent El fabricante del HLK-10M05T especifica una corriente de entrada
máxima de 0,2~A y recomienda un fusible de \textbf{1~A}/250~VCA de acción
retardada. El diseño usa \texttt{JFC2410-0200TS} de \textbf{0,2~A}. Con la carga
calculada, el fusible trabaja entre el 14~\pct{} y el 50~\pct{} de su corriente
nominal, lo que es correcto en régimen permanente, pero deja poco margen para la
corriente de irrupción.

\subsection{Coordinación del fusible con la irrupción}

\begin{samepage}
El \texttt{JFC2410-0200TS} tiene una resistencia en frío de 1\,646~m$\Omega$ y
un $I^2t$ de fusión de \textbf{0,0682~A$^2$s}. El HLK declara una corriente de
irrupción de hasta 10~A. Para un pulso exponencial de pico $I_0$ y constante de
tiempo $\tau$,
\begin{equation}
I^2t = \int_0^\infty \left(I_0 e^{-t/\tau}\right)^2 dt = \frac{I_0^2\,\tau}{2}.
\end{equation}
\end{samepage}

Con $I_0 = 10$~A, el $I^2t$ del pulso iguala al del fusible cuando
$\tau \approx 1{,}4$~ms. Para un pulso de 1~ms el consumo es 0,050~A$^2$s, es
decir \textbf{73~\pct{} del $I^2t$ de fusión en un solo arranque en frío}; para
2~ms lo supera. La holgura es insuficiente para un equipo que puede recibir
cortes repetidos.

\begin{table}[H]
\centering
\caption{Verificación del fusible \texttt{F300}--\texttt{F302}.}
\label{tab:fusible}
\small
\begin{tabularx}{\textwidth}{@{}l X l@{}}
\toprule
\textbf{Criterio} & \textbf{Resultado} & \textbf{Estado} \\
\midrule
Corriente permanente & 28--100~mA sobre 200~mA (14--50~\pct) & \ok \\
Caída de tensión en frío & 1,646~$\Omega \times 0{,}1$~A $= 165$~mV (0,14~\pct) & \ok \\
$I^2t$ de irrupción & 50--100~\pct{} del valor de fusión según $\tau$ & \revisar \\
Recomendación del fabricante & 1~A / 250~VCA retardado & \revisar \\
\bottomrule
\end{tabularx}
\end{table}

\noindent La coordinación también debe considerar los varistores
\texttt{MOV-14D471K} (470~V, 300~VCA máximos continuos) frente a la tensión de
fase-neutro de 120~V: trabajan al 40~\pct{} de su máximo continuo, lo que es
correcto, pero el fusible debe soportar la energía de una descarga sin abrirse
en un evento transitorio.

% ----------------------------------------------------------------------------
\section{Respaldo \emph{last gasp}}
\label{sec:respaldo}
% ----------------------------------------------------------------------------

\subsection{Energía del banco}

El banco son dos supercapacitores \texttt{HP-2R7-J306VYJ23} de 30~F y 2,7~V en
serie: $C_{eq} = 15$~F con ESR de 35~m$\Omega$ en continua cada uno
(70~m$\Omega$ el par). La tensión de corte no es la del banco sino la que exige
la cadena de salida. El buck del BOM solo puede reducir tensión, así que deja de
regular cuando su entrada cae por debajo de la salida más la caída de D201 y la
del interruptor interno:

\begin{equation}
V_{banco,min} = V_{out} + V_F(D201) + V_{drop}
= 3{,}3 + 0{,}42 + 0{,}06 \approx 3{,}78~\mathrm{V}.
\end{equation}

\noindent Esa es la diferencia conceptual con un respaldo sobre batería: la
energía útil no la fija el capacitor sino el convertidor. El \texttt{HX1002}
tiene capacidad de ciclo de trabajo del 100~\pct{}, lo que permite exprimir el
banco hasta ese punto; aun así, por debajo de 3,78~V la salida de 3,3~V empieza
a caer.

\begin{table}[H]
\centering
\caption{Energía útil y autonomía del banco.}
\label{tab:energia}
\small
\begin{tabularx}{\textwidth}{@{}X r r r r@{}}
\toprule
\textbf{Ventana de descarga} & \textbf{$E$ (J)} & \textbf{$t$ a 1~W} & \textbf{$t$ a 2~W} & \textbf{$t$ a 3,2~W} \\
\midrule
5,00 $\to$ 3,78~V (carga directa) & 80,3 & 80~s & 40~s & 25~s \\
4,58 $\to$ 3,78~V (carga tras D200) & 50,2 & 50~s & 25~s & 16~s \\
\bottomrule
\end{tabularx}
\end{table}

\noindent Dos observaciones sobre la tabla. La primera: el banco rinde entre
\textbf{50 y 80~J}, no los 24~J de una ventana conservadora hasta 4,2~V; la
tensión de corte real la pone el convertidor. La segunda, más importante: a
3,2~W el buck del BOM entrega 0,97~A, por encima de sus 0,7~A. \textbf{La carga
de cierre debe caber en 2,3~W} para ser compatible con el buck instalado; a ese
valor la autonomía es de 35~s con el banco cargado directamente. La limitación
del respaldo, hoy, no es la energía: es el regulador.

\subsection{Corriente y tiempo de carga}

La rama de carga del esquemático lleva \texttt{R207} de 3,9~$\Omega$ / 2~W en
serie. Desde un banco descargado:

\begin{table}[H]
\centering
\caption{Carga del banco a través de \texttt{R207}.}
\label{tab:carga}
\small
\begin{tabularx}{\textwidth}{@{}X r r r@{}}
\toprule
\textbf{Condición} & \textbf{Corriente inicial} & \textbf{Disipación en R207} & \textbf{Tiempo a 4,2~V} \\
\midrule
Banco vacío (0~V) & 1,28~A & 6,4~W & 107~s \\
Banco a 4,2~V & 0,21~A & 0,16~W & --- \\
\bottomrule
\end{tabularx}
\end{table}

\noindent \textbf{Observación.} Una resistencia de 2~W disipando 6,4~W durante
los primeros segundos de una carga desde cero queda fuera de su régimen. Si el
equipo puede arrancar con el banco descargado (por ejemplo, tras un fin de
semana sin alimentación), \texttt{R207} es un punto débil. Debe confirmarse en
Altium si hay otra resistencia en serie (el BOM tiene dos de 3,9~$\Omega$ /
2~W: R201 y R207) y cuál es la topología real de la rama.

\subsection{Autodescarga por las resistencias de balanceo}

El balanceo entre los dos capacitores en serie se hace con resistencias de
1~k$\Omega$ y 330~$\Omega$. Con 2~k$\Omega$ totales, la corriente de fuga del
banco es de unos 2,3~mA y la constante de autodescarga,
$\tau = 15~\mathrm{F}\times 2~\mathrm{k}\Omega \approx 8{,}3$~h. Un banco
cargado se vacía en menos de un día con el equipo apagado; el tiempo de recarga
(107~s) es aceptable, pero conviene que el firmware no intente un cierre
ordenado en el primer arranque tras un corte prolongado.

% ----------------------------------------------------------------------------
\subsection{¿Por qué supercapacitor y no batería?}
\label{sec:bateria}

La rúbrica pide justificar la química y la capacidad de la batería, si aplica.
Aquí \textbf{no aplica}: el producto se alimenta de la red y no está pensado
para operar sin ella. El banco no es una batería, es un depósito de energía para
un cierre ordenado. Vale la pena cuantificar la diferencia para que la decisión
quede documentada:

\begin{table}[H]
\centering
\caption{Equivalencia energética del banco frente a baterías comerciales.}
\label{tab:bateria}
\small
\begin{tabularx}{\textwidth}{@{}l r r X@{}}
\toprule
\textbf{Almacenamiento} & \textbf{Energía} & \textbf{Potencia} & \textbf{Comentario} \\
\midrule
2$\times$30~F (este diseño) & 6,6~mWh & 3,2~W & 1,5~mAh equivalentes ($Q=C\Delta V$) \\
CR2032 (primaria) & 660~mWh & $<10$~mW & 100 veces más energía, tres órdenes menos potencia \\
LiPo 100~mAh & 370~mWh & $\sim$1~W & Requiere protección de celda y carga CC/CV \\
\bottomrule
\end{tabularx}
\end{table}

\noindent El banco guarda \textbf{1,5~mAh} ($Q = C\,\Delta V = 15~\mathrm{F}
\times 0{,}36~\mathrm{V} = 5{,}4$~C) y \textbf{6,6~mWh}: cien veces menos energía
que una CR2032. Lo que ofrece a cambio es potencia y vida:

\begin{itemize}
  \item \textbf{Potencia.} Entrega 3,2~W sin protección ni electrónica de
        carga; una celda de litio de ese tamaño no lo hace sin un circuito de
        protección.
  \item \textbf{Vida cíclica.} El fabricante especifica 500\,000 ciclos entre
        tensión nominal y mitad de tensión, frente a los cientos de ciclos de
        una LiPo.
  \item \textbf{Temperatura.} $-40$ a $+70$~$^\circ$C, sin riesgo térmico ni de
        inflamación dentro del gabinete.
  \item \textbf{Carga.} 107~s desde vacío (Tabla~\ref{tab:carga}), sin
        protocolo CC/CV ni controlador dedicado.
\end{itemize}

\noindent\textbf{Limitación explícita.} El banco no da autonomía de horas. Si el
producto tuviera que operar sin red, la arquitectura cambiaría a una celda
LiFePO$_4$ con cargador y gestión, y el presupuesto debería rehacerse con la
energía como restricción principal en lugar de la potencia.

% ----------------------------------------------------------------------------
\section{Verificación térmica preliminar}
\label{sec:termica}
% ----------------------------------------------------------------------------

\begin{table}[H]
\centering
\caption{Disipación y aumento de temperatura estimado (25~$^\circ$C de ambiente,
$\theta_{JA}$ de hoja de datos o estimada).}
\label{tab:termica}
\small
\begin{tabularx}{\textwidth}{@{}X r r r l@{}}
\toprule
\textbf{Componente} & \textbf{$P_D$ típ.} & \textbf{$P_D$ máx.} & \textbf{$\Delta T$ máx.} & \textbf{Comentario} \\
\midrule
PS200 HLK-10M05T & 0,36~W & 1,00~W & --- & Módulo; límite de ambiente 60~$^\circ$C \\
U200 HX1002 (SOT-23-5) & 0,16~W & 0,48~W & $\sim$119~$^\circ$C & $\theta_{JA}\approx250$~$^\circ$C/W (estimada) \\
Alternativa AP63203 (TSOT-26) & 0,16~W & 0,48~W & 42~$^\circ$C & $\theta_{JA}=89$~$^\circ$C/W (hoja de datos) \\
D200 / D201 (SOD-323) & 0,12~W & 0,29~W & $\sim$102~$^\circ$C & $\theta_{JA}\approx350$~$^\circ$C/W (estimada) \\
U202 B0503S-1WR3 & 0,009~W & 0,012~W & --- & Módulo; despreciable \\
ESP32-S3-WROOM-1-N4 & 0,31~W & 1,20~W & --- & Pico de RF; requiere plano térmico \\
U500 RSM3485PCT & 0,10~W & 0,43~W & --- & Módulo aislado \\
\bottomrule
\end{tabularx}
\end{table}

\noindent \textbf{Lectura.} En operación típica ningún componente se acerca a su
límite. En el peor caso simultáneo, en cambio:

\begin{itemize}
  \item El \textbf{HX1002} en SOT-23-5 disiparía 0,48~W. Con 250~$^\circ$C/W y
        60~$^\circ$C de ambiente interior, la unión llegaría a
        $\sim$179~$^\circ$C, por encima de un límite típico de 125~$^\circ$C. Un
        buck en TSOT-26 con 89~$^\circ$C/W quedaría en 42~$^\circ$C de aumento:
        el reemplazo de U200 no debe resolverse solo por corriente, la
        disipación también obliga a cambiar de encapsulado.
  \item Los diodos de OR-ing en SOD-323 son el segundo punto caliente: 0,29~W en
        un encapsulado sin plano térmico. Recomendación: área de cobre generosa
        en el ánodo y verificación con el prototipo.
  \item El \textbf{HLK-10M05T} admite 60~$^\circ$C de ambiente. En un gabinete
        cerrado con 1,5~W de pérdidas internas, la temperatura interior puede
        superar ese valor en verano. Es el límite térmico del sistema, no de un
        componente aislado.
\end{itemize}

% ----------------------------------------------------------------------------
\section{Sensibilidad e incertidumbres}
\label{sec:sensibilidad}
% ----------------------------------------------------------------------------

Los dos parámetros que más mueven el resultado son la corriente del módulo
HaLow (sin dato de fabricante) y la capacidad real del buck.

\begin{table}[H]
\centering
\caption{Escenarios de carga del riel \texttt{DVCC\_3v3} y utilización del buck.}
\label{tab:escenarios}
\small
\begin{tabularx}{\textwidth}{@{}X r r r@{}}
\toprule
\textbf{Escenario} & \textbf{$I_{3,3}$ (mA)} & \textbf{HX1002 (700~mA)} & \textbf{Buck de 2~A} \\
\midrule
Optimista (HaLow 100~mA, SD 100~mA) & 710 & 101~\pct & 36~\pct \\
Base (Tabla~\ref{tab:dvcc}: HaLow 350~mA, SD 200~mA) & 1\,060 & 151~\pct & 53~\pct \\
Con reserva de diseño de 500~mA en el ESP32 & 1\,195 & 171~\pct & 60~\pct \\
HaLow en 500~mA (hipótesis de trabajo) & 1\,210 & 173~\pct & 61~\pct \\
\bottomrule
\end{tabularx}
\end{table}

\noindent No hay escenario en el que el \texttt{HX1002} tenga margen: incluso el
más favorable (HaLow y microSD a la mitad) lo deja al 101~\pct{} de su capacidad
nominal, y el caso base al 151~\pct{}. La disipación acompaña: a 710~mA el
encapsulado SOT-23-5 sube unos 65~$^\circ$C, de modo que con 40~$^\circ$C de
ambiente interior la unión ya roza los 105~$^\circ$C. Con un buck de 2~A la
utilización quedaría entre el 36~\pct{} y el 61~\pct{} en todos los casos.

\begin{table}[H]
\centering
\caption{Sensibilidad de la potencia de entrada al factor de potencia.}
\label{tab:pf}
\small
\begin{tabularx}{\textwidth}{@{}X r r r@{}}
\toprule
\textbf{PF} & \textbf{$I_{in}$ típico (mA)} & \textbf{$I_{in}$ peor caso (mA)} & \textbf{Con FEM (mA)} \\
\midrule
0,50 & 33,3 & 83,7 & 120,2 \\
0,60 & 27,8 & 69,7 & 100,1 \\
0,70 & 23,8 & 59,7 & 85,8 \\
\bottomrule
\end{tabularx}
\end{table}

\noindent Aun con PF~$=0{,}5$ y 100~V de entrada, la corriente de entrada no
supera los 120~mA con la rama FEM activa. El fusible de 0,2~A no es un problema
de régimen permanente; el riesgo está en la irrupción.

% ----------------------------------------------------------------------------
\section{Verificación contra el BOM}
\label{sec:discrepancias}
% ----------------------------------------------------------------------------

Esta sección confronta el BOM y el esquemático con la hoja de datos de cada
componente. No corrige nada: reporta y deja la decisión al equipo.

\begin{table}[H]
\centering
\caption{Verificación de los componentes del BOM contra su hoja de datos.}
\label{tab:discrepancias}
\small
\begin{tabularx}{\textwidth}{@{}>{\raggedright\arraybackslash}p{2.7cm}X>{\raggedright\arraybackslash}p{3.3cm}@{}}
\toprule
\textbf{Tema} & \textbf{Evidencia} & \textbf{Impacto} \\
\midrule
Buck U200 & El BOM lista \texttt{HX1002-GEC} con ``Output Current 1,3~A''. La hoja de datos del fabricante titula el dispositivo \emph{700~mA Synchronous Buck} y limita el pico del inductor a 0,9~A. & \critico{} -- La carga de pico (1\,060~mA) es el 151~\pct{} de 700~mA. Prevalece la hoja de datos. \\
\addlinespace
Carga del banco & El BOM y el esquemático resuelven la carga con \texttt{R207}/\texttt{R201} de 3,9~$\Omega$ / 2~W; no hay interruptor de corriente limitada. & \revisar{} -- La carga inicial disipa 6,4~W en una resistencia de 2~W (Tabla~\ref{tab:carga}). \\
\addlinespace
FGH100M & No figura en el BOM: se conecta por los headers 500/501. Su hoja de datos V1.0.0 (preliminar) no publica corrientes (\texttt{TBD}) y declara \texttt{VBAT} 3,0--3,6~V y \texttt{VDD\_FEM} 5~V típicos. & \revisar{} -- Los 350~mA son una asignación sin respaldo. La rama de 5~V no está en el presupuesto. \\
\addlinespace
Fusible F300 & La hoja de datos del \texttt{HLK-10M05T} recomienda 1~A retardado; el BOM trae \texttt{JFC2410-0200TS} de 0,2~A. & \revisar{} -- Ver Tabla~\ref{tab:fusible}. \\
\addlinespace
No están en el BOM & \texttt{AP63203}, \texttt{TPS2553}, \texttt{SY6280} y \texttt{HLK-5M03B} aparecen en hojas de datos y en documentación descriptiva, pero no son parte del diseño capturado. & Informativo: son alternativas evaluadas, no componentes de la placa. \\
\bottomrule
\end{tabularx}
\end{table}

\noindent \textbf{Nota de método.} Las tres primeras filas no son erratas
documentales: cambian el resultado del presupuesto. La última sirve para que la
revisión previa a fabricación se haga sobre una sola versión de la verdad: la
del BOM.

% ----------------------------------------------------------------------------
\section{Conclusiones y recomendaciones}
\label{sec:conclusiones}
% ----------------------------------------------------------------------------

\begin{table}[H]
\centering
\caption{Resumen ejecutivo del presupuesto de potencia.}
\label{tab:resumen}
\small
\begin{tabularx}{\textwidth}{@{}X r l@{}}
\toprule
\textbf{Indicador} & \textbf{Valor} & \textbf{Estado} \\
\midrule
Carga promedio de 3,3~V (perfil de 60~s) & 216~mA / 0,71~W & \ok \\
Carga de peor caso de 3,3~V & 1\,060~mA / 3,53~W & --- \\
Capacidad del buck del BOM (hoja de datos) & 700~mA & \critico \\
Capacidad del buck según el parámetro del BOM & 1\,300~mA & \revisar \\
Carga del riel aislado & 7,6--9,8~mA (3~\pct{} de 1~W) & \ok \\
Potencia de entrada (sin / con FEM) & 5,0~W / 7,2~W & \ok \\
Utilización del HLK-10M05T & 40--58~\pct & \ok \\
Energía útil del respaldo & 50--80~J & \ok \\
Autonomía del respaldo a 3,2~W & 16--25~s & \ok \\
Disipación del buck en el peor caso & 0,48~W & \revisar \\
Corriente de entrada frente al fusible & 14--50~\pct{} / irrupción al 73--100~\pct{} del $I^2t$ & \revisar \\
\bottomrule
\end{tabularx}
\end{table}

\noindent\textbf{Recomendaciones, en orden de prioridad:}
\begin{enumerate}
  \item \textbf{Resolver U200.} El BOM monta un \texttt{HX1002-GEC} de 700~mA y
        la carga de pico es 1,06~A. Hay que cambiar el buck por uno de al menos
        1,6~A (con margen) o reducir la carga, por ejemplo limitando el pico del
        HaLow por firmware. La decisión condiciona el resto del presupuesto.
  \item \textbf{Medir el FGH100M.} Su consumo es el mayor término no verificado
        del presupuesto y además tiene una rama de 5~V que hoy no se contabiliza.
  \item \textbf{Revisar la rama de carga del banco.} Verificar la topología real
        de \texttt{R201}/\texttt{R207} y su disipación en el arranque desde cero;
        si corresponde, volver al interruptor con límite programable.
  \item \textbf{Validar el fusible.} Medir la irrupción del HLK en frío y
        compararla con el $I^2t$ de 0,0682~A$^2$s; considerar 0,5--1~A
        retardado.
  \item \textbf{Verificar la térmica en gabinete.} El límite de 60~$^\circ$C del
        HLK y la unión del buck son los dos puntos que definen la temperatura
        interior admisible.
  \item \textbf{Actualizar la documentación descriptiva.} El \texttt{README},
        \texttt{AGENTS.md} y el diagrama de bloques describen otra revisión:
        citan un buck de 2~A, un interruptor de carga y un módulo AC/DC que no
        están en el BOM. Deben alinearse con el BOM y el esquemático.
\end{enumerate}

% ----------------------------------------------------------------------------
\section{Plan de verificación en banco}
\label{sec:verificacion}
% ----------------------------------------------------------------------------

\begin{table}[H]
\centering
\caption{Ensayos mínimos para cerrar el presupuesto con mediciones.}
\label{tab:ensayos}
\small
\begin{tabularx}{\textwidth}{@{}c X X@{}}
\toprule
\textbf{\#} & \textbf{Ensayo} & \textbf{Criterio de aceptación} \\
\midrule
1 & Consumo de \texttt{DVCC\_3v3} por estado: reposo, medición, TX Wi-Fi, escritura SD & Cada estado dentro del 15~\pct{} del valor de la Tabla~\ref{tab:dvcc} \\
2 & Consumo del FGH100M en TX a máxima potencia, en \texttt{VBAT} y \texttt{VDD\_FEM} & Dato para reemplazar la asignación de 350~mA \\
3 & Eficiencia del buck a 0,3 / 0,6 / 1,0~A & $\eta \ge 85$~\pct{} y $\Delta T < 60$~$^\circ$C \\
4 & Irrupción del HLK en frío (sonda de corriente, 10~ms) & $I^2t < 0{,}05$~A$^2$s para no consumir más del 70~\pct{} del fusible \\
5 & Autonomía real del respaldo con la carga de cierre & $\ge 5$~s con margen del 50~\pct{} \\
6 & Temperatura de U200, D200 y del HLK en cámara a 40~$^\circ$C & Unión $<110$~$^\circ$C; cápsula del HLK $<60$~$^\circ$C \\
7 & Autodescarga del banco a las 24~h & Consistente con $\tau\approx8{,}3$~h \\
\bottomrule
\end{tabularx}
\end{table}

% ----------------------------------------------------------------------------
\section*{Referencias}
\addcontentsline{toc}{section}{Referencias}
% ----------------------------------------------------------------------------

\begin{enumerate}[label={[\arabic*]},leftmargin=2.2em]
  \item Hi-Link. \emph{10W Ultra-Compact Power Module Series}
        (\texttt{HLK-10M05\_Datasheet.pdf}), V1.0, dic. 2024.
  \item HEXIN. \emph{HX1002 -- 700~mA Synchronous Buck DC/DC Converter}
        (\texttt{HX1002-GEC\_Datasheet.pdf}).
  \item Diodes/AP. \emph{AP63203 -- 2~A Synchronous Buck Converter}
        (\texttt{AP63203WU-7\_Datasheet.pdf}).
  \item Texas Instruments. \emph{ISO7741 -- Quad-Channel Digital Isolator}
        (\texttt{ISO7741DWR\_Datasheet.pdf}), SLLSEP4G, feb. 2020.
  \item Renergy. \emph{RN8302B Polyphase Energy Metering IC}
        (\texttt{RN8302B\_Datasheet\_v3.pdf}).
  \item Espressif. \emph{ESP32-S3-WROOM-1 Datasheet} v1.8
        (\texttt{ESP32-S3-WROOM\_Module.pdf}).
  \item Visom. \emph{RSM3485PCT -- RS-485 aislado}
        (\texttt{RSM3485PCT\_Datasheet.pdf}).
  \item EVISUN. \emph{B0503S-1WR3 -- DC-DC aislado 1~W}
        (\texttt{EVISUN B0503S-1WR3.pdf}).
  \item Texas Instruments. \emph{TPS2553 -- Adjustable Current-Limited Power
        Switch} (\texttt{TI TPS2553DRVR.pdf}).
  \item Quectel. \emph{FGH100M-H Wi-Fi HaLow Module Specification}, V1.0.0
        preliminar, oct. 2024.
  \item KAMCAP. \emph{HP-2R7-J306VYJ23 -- Supercapacitor 30~F / 2,7~V}
        (\texttt{Supercapacitor\_Datasheet.pdf}).
  \item JDT Fuse. \emph{JFC2410TS Time-Lag Fuse}
        (\texttt{JFC2410-0200TS\_Datasheet.pdf}), rev. 11/05/25.
  \item MSKSEMI. \emph{MS2A40LWS -- Schottky 2~A / 40~V}
        (\texttt{MS2A40LWS\_Datasheet.pdf}).
  \item Texas Instruments. \emph{TIDA-010244: Three-Phase Shunt-Based Energy
        Metrology Reference Design} (\texttt{TI\_DUF27A\_Application\_Note.pdf}).
  \item Equipo del proyecto. Esquemático
        (\path{PCB/SCH/Analizador_trifasico_IoT_pcb.pdf}) y BOM
        (\path{PCB/BOOM/Analizador_trifasico_IoT_pcb.xlsx}), 15/11/2026.
\end{enumerate}

% ----------------------------------------------------------------------------
\appendix
\section{Parámetros extraídos por componente}
% ----------------------------------------------------------------------------

\begin{table}[H]
\centering
\caption{Parámetros de hoja de datos usados en los cálculos.}
\label{tab:parametros}
\small
\begin{tabularx}{\textwidth}{@{}l X@{}}
\toprule
\textbf{Componente} & \textbf{Parámetros relevantes} \\
\midrule
HLK-10M05T & 10~W; 5~V/2~A (2,6~A corto); 85--265~VCA; $\eta \ge 80$~\pct{} a 115~V; sin carga $<0{,}1$~W; rizado $\le100$~mV; aislamiento 3000~VCA; ambiente $-25$ a $+60$~$^\circ$C; irrupción $\le10$~A \\
HX1002 & 0,7~A; 2,2--5,5~V; $\eta$ hasta 95~\pct; $I_Q=200~\mu$A; límite de pico 0,9~A; $R_{DS(on)}$ 0,30/0,39~$\Omega$; SOT-23-5L \\
AP63203 (alternativa) & 2~A; 3,8--32~V; $I_Q=22~\mu$A; hasta 88~\pct{} en carga ligera; $\theta_{JA}=89$~$^\circ$C/W (TSOT-26) \\
ISO7741 & 4 canales (3/1); 5000~V$_{RMS}$; 3,3~V: 3,2~mA lado 1 y 2,7~mA lado 2 a 1~Mbps \\
RN8302B & 3,0--3,6~V; $I_{DD}=4{,}4$~mA (EMM); 4,7~mA máx.; ADC 4,1~mA; NVM 2,3~mA \\
RSM3485PCT & 3,15--3,45~V; 30~mA típ., 130~mA máx.; 2500~VCC; 9600~bps; 64 nodos; CON 5~mA \\
B0503S-1WR3 & 5~V$\to$3,3~V; 1~W; 303~mA máx.; $\eta=74$~\pct; 1500~VCC \\

FGH100M-H & \texttt{VBAT} 3,0--3,6~V; \texttt{VDD\_FEM} 3,0--5,25~V (típ.\ 5~V); 27~dBm máx.; 32,5~Mbps; corrientes \texttt{TBD} \\
HP-2R7-J306VYJ23 & 30~F; 2,7~V; ESR 22~m$\Omega$ (AC) / 35~m$\Omega$ (DC); fuga 0,09~mA; 6,7~A nominal \\
MS2A40LWS & 40~V; 2~A; $V_F=0{,}37$~V a 1~A; 0,50~V a 2~A; SOD-323 \\
JFC2410-0200TS & 0,2~A; 250~V; 1,646~$\Omega$; $I^2t=0{,}0682$~A$^2$s; 200~\pct{} en 60~s máx. \\
MOV-14D471K & 470~V; 300~VCA continuos; 4,5~kA \\
\bottomrule
\end{tabularx}
\end{table}

% ----------------------------------------------------------------------------
\section{Expresiones y valores intermedios}
% ----------------------------------------------------------------------------

\noindent\textbf{Carga del riel digital.}
\begin{align}
I_{3,3,max} &= 355 + 10 + 200 + 130 + 350 + 4{,}6 + 10 \nonumber\\
            &= 1\,059{,}6~\mathrm{mA} \approx 1{,}06~\mathrm{A}.
\end{align}

\noindent\textbf{Potencia del riel digital.}
\begin{align}
P_{3,3,max} &= 3{,}3~\mathrm{V}\times1{,}0596~\mathrm{A} = 3{,}497~\mathrm{W}.
\end{align}

\noindent\textbf{Entrada del buck.}
\begin{align}
P_{5,buck} &= \frac{3{,}497}{0{,}88} = 3{,}974~\mathrm{W},
& I_{5,buck} &= 794{,}7~\mathrm{mA}.
\end{align}

\noindent\textbf{Entrada del DC-DC aislado.}
\begin{align}
P_{5,iso} &= \frac{3{,}3\times0{,}0098}{0{,}74} = 43{,}7~\mathrm{mW},
& I_{5,iso} &= 8{,}7~\mathrm{mA}.
\end{align}

\noindent\textbf{Entrada de CA (peor caso, sin FEM).}
\begin{align}
P_{AC} &= \frac{5\times0{,}8034}{0{,}80} = 5{,}02~\mathrm{W},
& I_{AC} &= \frac{5{,}02}{120\times0{,}60} = 69{,}7~\mathrm{mA}.
\end{align}

\noindent\textbf{Energía del respaldo (5,00~V a 3,78~V).}
\begin{align}
E &= \tfrac{1}{2}\times15\times\left(5{,}00^2-3{,}78^2\right) = 80{,}3~\mathrm{J},
& t_{3,2W} &= \frac{80{,}3}{3{,}2} = 25~\mathrm{s}.
\end{align}

\noindent\textbf{Potencia promedio del riel de 3,3~V.}
\begin{align}
P_{3,3,prom} &= 3{,}3\times(0{,}0902+0{,}0052+0{,}0400+0{,}0599+0{,}0032+0{,}0100+0{,}0076) \nonumber\\
&= 3{,}3\times0{,}2161 = 0{,}713~\mathrm{W}.
\end{align}

\end{document}
